%\subsection{Rhythmic analysis}\label{ch02_sb:sec:rhythmic-analysis}
%The rhythmic structure in IAM is based on the tāla framework, which is a hierarchical and fixed-duration metric that provides clues for rendition, repetition of melodic patterns, and more. The said hierarchical form of a tāla is targeted in~\citep{ajay_rhythm_hierarchical}, using a supervised approach on top of spectral hand-crafted features that capture tempo and onsets in the signal.

%Bayesian models are found to be the most capable of capture the rhythmic cycles in music signals, first applied in the form of particle filters for an efficient inference~\citep{ajay-rhythm_particle}, showing competetitive results on Carnatic recordings. The said approach is further extended to a generalized model to capture long cycles~\citep{ajay-rhythm_bayesian}, being this very convenient for rhythm tracking in Hindustani Music.
%
%Relevant research has been also carried out aiming at characterizing percussion patterns in the most common rhythmic instruments in IAM~\citep{gupta-syllabic}.


%\subsection{Structural analysis}\label{ch02_sb:sec:structural-analysis}
%Carnatic and Hindustani performances have a very tradition-specific organization in sections, including the main core of the rāga, improvisation, instrument solos, and more. There is not a single specific or standard form in the two traditions: structure is very much dependant on the type of composition. Despite that, sections in IAM performances are notably distinguishable in terms of melody, rhythm, and instrumentation. For that reason, the structural analysis of Carnatic and Hindustani music can be approached using whether rhythmic or melodic features.

%\citep{verma-structural} proposes to use tempo cues combined with additional hand-crafted features that aim at capturing tradition-relevant information. On the other hand, \citep{ranjani-structural} quantize the pitch contours, identify repeated note sequences or patterns, and then compute posterior probabilities to identify the different sections. Alternatively, \citep{ganguli-structures} builds on top of hand-crafted features of pitch time-series.


%\subsection{Timbre analysis}\label{ch02_sb:sec:timbre-analysis}
%Not a lot of research has been carried out towards instrument identification or classification in an IAM context, since the instrument arrangement and performance setup is normally consistent along distributed recordings, while isolated instrument audio signals are scarce. However, percussion instruments in IAM have, in themselves, a list of particular strokes which are very relevant to classify, whether using non-negative matrix factorization for feature processing~\citep{anantapadmanabhan-mridangam_1,anantapadmanabhan-mridangam_2}, or leveraging from DL methods for drums~\citep{rohit-tabla}.
%being the mridangam the typical case in Carnatic~\citep{mridangam1, mridangam2}, and tabla in Hindustani~\citep{rohit-tabla}.
